Для $x=(x_1,x_2)$ чебышёвский базис состоит из произведений $T_j(x_1/2)T_l(x_2/2)$. Для Fourier-SNO используются вещественные косинусы $\cos(j\pi(x_1+2)/4)\cos(l\pi(x_2+2)/4)$, $0\le j,l<16$. Этот выбор соответствует чётному продолжению по каждой координате с периодом 8 и допускает различные значения на противоположных сторонах квадрата. Обе модели используют 256 коэффициентов на компоненту. Матрица $V$ содержит значения базиса в узлах; коэффициенты проекции решения определяются из $(V^\T M V)c=V^\T M f$. Семь параметрических признаков линейно переводятся в таблицу $16\times16$ с 24 каналами. В четырёх остаточных блоках обучаемые преобразования действуют по двум осям мод и по каналам, после чего применяется GELU. Такое построение адаптирует коэффициентный слой из приложения~C работы~\cite{fo} к непериодической задаче и треугольной сетке.
